\documentclass[11pt]{article}
\usepackage[T1]{fontenc}
\usepackage{textcomp}
\usepackage[margin=0.75in]{geometry}
\usepackage{amsmath,amssymb,amsthm}
\usepackage{array,booktabs}
\usepackage{hyperref}
\setlength{\parindent}{0pt}
\newtheorem{theorem}{Theorem}
\theoremstyle{definition}
\newtheorem{definition}{Definition}
\title{Flow Proof Artifact: prop-03-circumscribe-equiangular-triangle}
\author{Generated by \texttt{flow doc proof}}
\date{Flow Proof Book}
\begin{document}
\maketitle
About a given circle to circumscribe a triangle equiangular to a given triangle.\\[0.5em]
\textbf{Source.} euclid — Elements, Book IV, Proposition 3\\[0.5em]
\bigskip
\noindent{\large \textbf{Derived fact 1.} \textit{About a given circle to circumscribe a triangle equiangular to a given triangle} \hfill the Euclidean plane \cdot Euclid Book IV \cdot Proposition 3: about a circle to circumscribe a triangle equiangular to a given triangle}\par
\smallskip\noindent\textit{Coordinate: the Euclidean plane \cdot Euclid Book IV \cdot Proposition 3: about a circle to circumscribe a triangle equiangular to a given triangle}\par
\smallskip\noindent\textit{Source: euclid — Elements, Book IV, Proposition 3}\par
\smallskip\noindent\textit{Built on: proposition 2: in a circle to inscribe a triangle equiangular to a given triangle, for Euclid Book IV on the Euclidean plane, proposition 19: the tangent is perpendicular to the diameter at the point of contact, for Euclid Book III on the Euclidean plane}\par
\medskip
\noindent\textbf{Goal.} About a given circle to circumscribe a triangle equiangular to a given triangle\par
\noindent$\displaystyle \text{circumscribed triangle equiangular}$\par
\medskip
\noindent
\begin{tabular}{@{} >{\raggedright\arraybackslash}p{0.06\textwidth} >{\raggedright\arraybackslash}p{0.47\textwidth} >{\raggedleft\arraybackslash}p{0.41\textwidth} @{}}
\textbf{\#} & \textbf{Proof} & \textbf{Mathematics} \\
\midrule
\textbf{1.} & We prove that proposition 3: about a circle to circumscribe a triangle equiangular to a given triangle for Euclid Book IV the Euclidean plane. & \\[0.45em]
\textbf{2.} & Let circle = given\_circle. & \\[0.45em]
\textbf{3.} & Let ABC = given\_triangle. & \\[0.45em]
\textbf{4.} & Let PQR = circumscribed\_triangle. & \\[0.45em]
\textbf{5.} & From step 2, step 3, and step 4, we can deduce that angle PQR equals angle ABC. & $\displaystyle \angle PQR = \angle ABC$ \\[0.45em]
\textbf{6.} & We invoke the derived fact governing Euclid Book IV the Euclidean plane: proposition 2: in a circle to inscribe a triangle equiangular to a given triangle, for Euclid Book IV on the Euclidean plane. & \\[0.45em]
\textbf{7.} & We invoke the derived fact governing Euclid Book III the Euclidean plane: proposition 19: the tangent is perpendicular to the diameter at the point of contact, for Euclid Book III on the Euclidean plane. & \\[0.45em]
\textbf{8.} & From step 6 and step 7, this implies circumscribed triangle equiangular. Hence proven. & $\displaystyle \text{circumscribed triangle equiangular}$ \\[0.45em]
\bottomrule
\end{tabular}
\label{thm:Geometry-Euclid Book IV-proposition_3_about_a_circle_to_circumscribe_a_triangle_equiangular_to_a_given_triangle}
\par\medskip\hrule
\end{document}
